% By default, the contents of the back side of the final sheet is
% printed upside down. The option notumble suppresses that. Doing so
% may be necessary to suit the behavior of certain printing engines.
% Specifying [notumble] may also be useful during the writing of a
% document, to enable proof-reading on the screen.
\documentclass[notumble]{leaflet}
%\documentclass{leaflet}


\usepackage[english]{babel}
\usepackage[utf8]{inputenc}
\usepackage{lmodern}
\usepackage[T1]{fontenc}

\usepackage[default,scale=1]{opensans}
% \renewcommand{\seriesdefault}{cl}
\renewcommand{\bfdefault}{sb}
% \renewcommand{\ttdefault}{fos}
\renewcommand\ttfamily{%
  \fontfamily{\ttdefault}%
  % \fontseries{lc}%
  \fontshape{\shapedefault}%
  \selectfont}

\usepackage{enumitem}
\setlist{
    topsep=0pt,
    partopsep=0pt,
    parsep=1ex,
    nolistsep,
    itemsep=.1ex
}
\setlength{\parskip}{1ex}

% \newcommand{\myparagraph}[1]{\vspace{.6em}\textsc{#1}}
\newcommand{\myparagraph}[1]{#1---}

\usepackage[table,svgnames]{xcolor}
\usepackage{graphicx}
\usepackage{tabularx,booktabs}
\usepackage{alltt,textcomp}
\usepackage{url}
\urlstyle{sf}
\usepackage{ifthen}
\usepackage{xspace}

\newcommand{\eg}{\emph{e.g.,}\xspace}
\newcommand{\ie}{\emph{i.e.,}\xspace}
\newcommand{\etal}{\emph{et al.,}\xspace}
\newcommand{\ct}[1]{{\textsf{#1}}\xspace}

\usepackage{microtype}

\definecolor{light-gray}{rgb}{0.8,0.8,0.8}
\definecolor{comment}{rgb}{0.35,0.35,0.35}
\definecolor{string}{rgb}{0.5,0,0.5}
\definecolor{darkRed}{rgb}{0.5,0,0}
\definecolor{darkBlue}{rgb}{0,0,0.3}
\definecolor{darkCyan}{rgb}{0,0.3,0.3}
\definecolor{link}{rgb}{0.16,0.59,0.84}

\makeatletter
\newsavebox{\code@box}
 \newenvironment{displaycode}{%
     \par
     % \begin{lrbox}{\code@box}%
     \hspace{1.5em}\begin{minipage}{\linewidth}
       \begin{alltt}\small}{
       \end{alltt}
     \end{minipage}
     % \end{lrbox}
     % \usebox{\code@box}%
     \par}
\makeatother


% Use of english prevents babel from adding spaces before $:
\newcommand{\code}[1]{\foreignlanguage{english}{\texttt{#1}}}

\graphicspath{{./figures/}}

\title{%
  % \fontseries{cl}\selectfont\huge
  % Cheat sheet\\
  \vspace*{0.5\textwidth}%
  \normalfont
  \includegraphics[width=\linewidth]{logo.pdf}
  \\[1.8\baselineskip]%
  \fontseries{cl}\selectfont\Huge%
  The Living\\ Language and Environment
  % a clean, innovative, open-source Smalltalk-inspired environment\\
  \vfill
  \url{http://www.pharo.org}}

\date{}

\setmargins{8mm}{8mm}{8mm}{8mm}

\CutLine*{1} \CutLine*{2} \CutLine*{3} \CutLine*{4} \CutLine*{5}
\CutLine*{6}

\usepackage[colorlinks=true, urlcolor=link, pdfborder={0 0 0}]{hyperref}
\begin{document}
\maketitle
\thispagestyle{empty}

\pagebreak{}

Pharo's goal is to deliver a clean, innovative, free open-source
living object-oriented language and environment. By providing a stable
and small core system, excellent developer tools, and maintained
releases, Pharo is an attractive platform to build and deploy mission
critical applications. Pharo focuses on developing a flexible and
advanced infrastructure that can be used by SMEs (Small and Medium
Enterprises) as well as researchers.

\section{Pharo as an ecosystem}

The goal of Pharo is to foster a healthy ecosystem of both private and
commercial contributors who advance and maintain the core system and
its external packages. In addition, Pharo wants to support a community
of researchers in the domain of reengineering, language design,
robotics and web development.

\section{Pharo’ forge}

Pharo is not only a core set of packages but more than 300 projects developed by external contributors and we should build an infrastructure to control the quality of the production.

\begin{center}
\url{http://www.smalltalkhub.com}
\end{center}


It contains around 900 projects, 700 users and 90000 packages.

\section{Books}

\begin{itemize}
\item Pharo by Example 2009 - (translated to french, spanish japanese)
\item Deep into Pharo (2013)
\end{itemize}

\pagebreak{}

\section{Pharo for teaching}

Universities use Pharo to teach:

\begin{itemize}
\item Uni. of Buenos Aires
\item Uni. of Bern
\item Uni. of Brussels
\item Ecole des Mines de Douai
\item Uni. de Savoie
\item Ivan Franko Nat. Uni. of Lviv
\item Czech Technical Uni.
\item CULS Prague
\item Uni. of Quilmes
\item Uni. of La Plata
\item Northern Michigan Uni.
\item Uni. Technologica Nacional (UTN)
\item Uni. Catholic of Argentina
\item Uni. of Santiago
\item Uni. Policnica de Catalunya
\item Uni. of Saarland
\item Uni. of San Martin
\item Uni. of Maroua
\item Uni. de Bretagne Occidentale
\item Uni. de Tomsk 
\end{itemize}

\section{Pharo as a platform for research}

Some research groups use Pharo:

\begin{itemize}
\item LAFHIS, AR
\item SCG, CH
\item CEA-List, FR
\item CAR-EMD, FR
\item RMOD-Inria, FR
\item UMMISCO-IRD, FR
\item Reveal, CH
\item LYSIC, FR
\item IGEVET, AR
\item Pleiad, CL
\end{itemize}

\pagebreak{}

\section{Pharo as a platform for business}

Companies use Pharo:

\begin{itemize}
\item \url{http://www.romaxtech.com}
\item \url{http://www.netstyle.ch}
\item \url{http://www.pinesoft.uk}
\item \url{http://www.cmsbox.com}
\item \url{http://smallworks.com.ar}
\item \url{http://www.2denker.de}
\item \url{http://www.inceptive.be}
\item \url{http://gemtalksystems.com}
\item \url{http://www.nidea.eu}
\item \url{http://www.doslovno.info}
\item \url{http://www.10pines.com}
\item \url{http://www.agilitic.com}
\item \url{http://www.ibizlog.com}
\item \url{http://www.ardishealth.com}
\item \url{http://www.airflowing.com}
\item \url{http://www.majcon.de}
\item \url{http://www.mindclue.ch}
\item \url{http://www.synectique.eu}
\item \url{http://www.objectprofile.com}
\item \url{http://www.objectfusion.fr}
\item \url{http://www.square-i.net}
\item \url{http://yesplan.be}
\item \url{http://debrispublishing.com}
\item \url{http://www.beta9.be}
\item \url{http://delware.nl}
\end{itemize}

Other projects cannot be publicly revealed: For example, Pharo is used to deploy software to more than 6000 users in a large aircraft company.

\section{Consortium}

Companies support Pharo via a consortium:
\begin{center}
  \url{http://consortium.pharo.org}
\end{center}


Advanced consultants are available:
\begin{center}
  \url{http://consultants.pharo.org}
\end{center}


\end{document}

%%% Local Variables:
%%% mode: latex
%%% TeX-master: t
%%% End:
